\paragraph{Registration costs and the composition of participation.}
A first strand studies how electoral rules alter the cost of becoming, or remaining, an eligible voter. Classic evidence from the United States shows that registration requirements depress turnout, especially when deadlines are early, offices are inconvenient, or reregistration is required after moving \cite{Rosenstone1978TheEffectRegistra,Ansolabehere2005TheIntroducVoter}. Subsequent work complicates the distributional claim: education gaps need not always widen mechanically when registration rules change \cite{Nagler1991TheEffectRegistra}, while reforms that lower administrative frictions, such as preregistration and automatic reregistration, can raise participation among targeted groups \cite{Holbein2014MakingYoungVoters,Kim2022AutomatiVoterReregist}. Comparative work similarly emphasizes that turnout reflects both rules and mobilization incentives \cite{cox2015,frankMartinezComa2023,gaeblerEtAl2020}. The voter-ID literature pushes this argument toward institutional design that is explicitly justified by electoral integrity: strict ID rules can screen out voters without compliant documentation, with especially large burdens for racial and ethnic minorities and durable effects even when rules are later relaxed \cite{Fraga2021WhoVoterLaws,Hajnal2017VoterIdentifiLaws,Grimmer2021TheDurableDifferen}. Other evidence bounds the aggregate consequences more tightly, showing that in Florida and Michigan only a very small share of ballots were cast without ID and that even extreme disenfranchisement assumptions would rarely change election outcomes \cite{Hoekstra2019Strict}. This paper studies a related but distinct reform: biometric voter registration does not merely require identification at the polling place, but changes the administrative stock of registered voters itself, making it possible to measure both aggregate registry exit and compositional change.

\paragraph{Election technology and biometric state capacity.}
A second strand asks whether election technology expands effective access or instead substitutes one form of error for another. In Brazil, electronic voting reduced residual votes and increased the political weight of less-educated voters, with downstream effects on representation and infant health \cite{Fujiwara2015VotingTechnoloPolitica}. Yet the same broad technological modernization agenda can introduce new mistakes: electronic voting also increased party-label votes in ways consistent with ballot-interface confusion \cite{Zucco2016TradingOldErrors}. The closest existing study of Brazilian biometric registration finds that biometric recadastramento reduced the number of registered voters but did not significantly reduce actual turnout, interpreting the gap partly as the removal of invalid registrations \cite{forquesatoRodrigues2023}. Outside electoral administration, biometric monitoring has been studied as a state-capacity tool: biometric smartcards in India reduced leakage and improved payment delivery without clear evidence of beneficiary exclusion, while biometric monitoring in health care improved adherence and data accuracy \cite{Muralidharan2016BuildingStateCapacity,Bossuroy2024BiometriMonitoriService}. The unresolved issue is whether the integrity gains from biometric administration come with politically meaningful changes in who remains registered. This paper addresses that gap by decomposing registry change into exit and education relabeling at the municipality-cohort level.

\paragraph{Unequal turnout, representation, and welfare.}
A third branch links participation to representation and welfare. If turnout is uneven, policy need not reflect the preferences of the full eligible population: in city politics, higher turnout among underrepresented groups can shift descriptive and substantive representation \cite{Hajnal2005WhereTurnoutMatters}. Political-economy models give the mechanism theoretical bite: voting costs, candidate entry, and the location of the decisive voter shape equilibrium policy, redistribution, and the welfare consequences of participation rules \cite{besleyCoate1997,Meltzer1981RationalTheoryThe,Borgers2004CostlyVoting,Krasa2008MandatorVotingBetter}. A related public-finance literature shows why administrative restrictions can improve targeting in principle while also screening out real beneficiaries in practice \cite{Nichols1982TargetinTransferThrough,Deshpande2019WhoScreenedOut}. Recent evidence from Brazil reinforces the stakes: restrictions on fraudulent registration can alter political competition, candidate selection, public goods, and social outcomes \cite{Karim2025VoterBuyingPolitici}. The contribution here is to bring these welfare and representation concerns to a measured electoral-administration shock: rather than assuming that biometric recadastramento only removes ghosts or only burdens real voters, the paper estimates the size and composition of the affected electorate and uses those estimates in a welfare framework.

\paragraph{Electoral integrity, administration, and institutional trust.}
A fourth strand studies how citizens interpret electoral administration itself. Electoral reforms may increase procedural integrity, but they can also change beliefs about the fairness or competence of institutions. Mobilization around electoral procedures in Kenya increased turnout but lowered trust in the electoral commission, particularly among politically disappointed voters and in violent constituencies \cite{Marx2021VoterMobilisaAnd}. International and domestic actors also value electoral process independently of outcomes, which is why process integrity can become politically salient even when candidate stakes are high \cite{Bubeck2017ProcessCandidatThe}. Evidence from local election administration suggests that partisan administrators do not necessarily shift vote shares in favor of their party, but the question of perceived neutrality remains central to democratic legitimacy \cite{Ferrer2023HowPartisanLocal}. More broadly, electoral accountability can discipline corruption when citizens can observe and sanction incumbents \cite{ferrazFinan2011}. Biometric voter registration sits at the intersection of these concerns: it is advertised as an anti-fraud reform, but it may also make registration loss more visible, create confusion, or change who is represented in the electorate. The LAPOP component of this paper therefore treats trust in institutions and democratic attitudes as outcomes, not merely as background conditions.

\paragraph{Gap and contribution.}
Taken together, existing work shows that administrative frictions affect turnout, that election technologies can enfranchise some voters while generating new errors, that unequal participation matters for representation and welfare, and that procedural reforms can reshape trust. What remains less well understood is how a large biometric registration reform changes the electorate before votes are cast, especially outside the U.S. voter-ID setting and in a context where biometric registration is rolled out gradually across municipalities. This paper fills that gap by combining Brazilian administrative data with survey evidence: it estimates municipality-level registry exit, education relabeling, downstream voting patterns, and institutional trust, while distinguishing strict and hybrid implementation regimes. The result is a bridge between the voter-ID literature's concern with burdens, the state-capacity literature's focus on administrative integrity, and political-economy models of representation and welfare.
